\bookpreface{Prefácio}

Esse livro surgiu do jeito que eu ensinava ML. Como vim da física, minha maneira de pensar em machine learning é bem diferente da de um cientista da computação médio ou de um estatístico médio. Eu nunca vi outro livro abordar o assunto desse jeito.

Nada aqui é particularmente novo, tirando o caminho que eu faço pra chegar aos resultados. Isso vem muito da minha formação em física, especialmente em física teórica. Vem também do fato de eu ser fã das \emph{Feynman Lectures on Physics}, que mudaram minha maneira de pensar sobre a física em geral.

Este livro não foi pensado pra quem está começando na área. Acho que quem tentasse aprender por aqui ia ter dificuldade, porque deixei de lado muita coisa clássica e boa parte da organização que costuma orientar um primeiro curso. Mas, se essa é a sua segunda ou terceira passada pelo assunto, você pode descobrir outra maneira de ver esse conteúdo e entender \emph{como} ele é usado na prática.

Os primeiros capítulos podem ser lidos lado a lado com \href{https://amlbook.com/}{\emph{Learning from Data}}, de Yaser S. Abu-Mostafa, Malik Magdon-Ismail e Hsuan-Tien Lin. Ali você vai encontrar um curso mais tradicional, que ensina os fundamentos matemáticos da área.

Particularmente, gosto bastante da distinção que fiz entre \emph{generative learning} e \emph{discriminative learning}. A maioria dos livros menciona as duas abordagens, mas acho que muita gente acaba a leitura sem ter internalizado a diferença entre os dois caminhos.
